\documentclass[11pt]{article}
\usepackage{amsmath,amssymb,amsthm}
\usepackage[margin=1in]{geometry}

\newtheorem{theorem}{Theorem}
\newtheorem{lemma}[theorem]{Lemma}
\newtheorem{proposition}[theorem]{Proposition}

\newcommand{\E}{\mathbb E}

\title{Solution to Ramanujan Challenge Problem 2.2:\\
Euler's Constant $\gamma$ as an Ap\'ery Limit}
\author{Xiang Huang}
\date{August 2026}

\begin{document}
\maketitle

\begin{abstract}
We prove unconditionally that the quotient defined by the four-term recurrence
in Problem~2.2 converges to Euler's constant.  The challenge recurrence is an
exact first-order Ore transform of a recurrence with positive hypergeometric
solutions.  For the latter, the quotient is a weighted average of harmonic
numbers.  A finite birth--death (Stein) identity gives a quantitative second
moment bound at the cubic saddle $k^3=(n-k)^2$; elementary harmonic-number
bounds then imply convergence to $\gamma$.  Thus no asymptotic theorem for a
resonant recurrence and no cited analytic limit is required.
\end{abstract}

\section{The challenge recurrence}

After shifting the three negative initial indices, write the challenge
solutions as $P_n=p_{n-3}$ and $Q_n=q_{n-3}$.  Thus
\[
 (P_0,P_1,P_2)=(0,7,179),\qquad
 (Q_0,Q_1,Q_2)=(1,12,306).
\]
They satisfy
\begin{equation}\label{eq:challenge}
 c_0(n)U_{n+3}+c_1(n)U_{n+2}+c_2(n)U_{n+1}+c_3(n)U_n=0,
\end{equation}
where
\begin{align*}
c_0(n)&=-8n^3-51n^2-105n-68,\\
c_1(n)&=24n^5+337n^4+1833n^3+4818n^2+6092n+2928,\\
c_2(n)&=-(n+2)(n+3)
 (24n^5+273n^4+1150n^3+2154n^2+1635n+268),\\
c_3(n)&=(n+1)(n+2)^4(n+3)(8n^3+75n^2+231n+232).
\end{align*}

We will prove that $P_n/Q_n\to\gamma$.  Shifting the index back then gives
the limit in the statement of the problem.

\section{A positive hypergeometric model}

Let $A_n,B_n$ be the two solutions of
\begin{align}
 &(n+3)^2(8n+11)(8n+19)R_{n+3}\nonumber\\
 &\quad -(n+3)(8n+11)(24n^2+145n+215)R_{n+2}\nonumber\\
 &\quad +(8n+27)(24n^3+105n^2+124n+25)R_{n+1}\nonumber\\
 &\quad -(n+2)^2(8n+19)(8n+27)R_n=0,                 \label{eq:rivoal}
\end{align}
with
\[
 (A_0,A_1,A_2)=(-1,4,77/4),\qquad
 (B_0,B_1,B_2)=(1,7,65/2).
\]
Put $H_m=\sum_{j=1}^m1/j$, with $H_0=0$, and
\[
 w_{n,k}=\frac{\binom nk^2}{k!}\quad(0\le k\le n).
\]

\begin{proposition}[finite-sum formulas]\label{prop:sums}
For every $n\ge0$,
\begin{align}
B_n&=\sum_{k=0}^n(2n+k+1)w_{n,k},                                      \label{eq:Bsum}\\
A_n&=\sum_{k=0}^n w_{n,k}
 \bigl((2n+k+1)(3H_k-2H_{n-k})-1\bigr).                               \label{eq:Asum}
\end{align}
In particular, $B_n>0$.
\end{proposition}

\begin{proof}
The adjacent quotient
\[
 \frac{w_{n,k+1}}{w_{n,k}}=\frac{(n-k)^2}{(k+1)^3}
\]
makes both summands hypergeometric (with the harmonic part handled by
$H_{k+1}-H_k=1/(k+1)$).  Substitution in~\eqref{eq:rivoal}, followed by
clearing denominators, is a telescoping identity in $k$; the boundary terms
vanish at $k=0$ and $k=n+1$.  The resulting sums have the displayed initial
triples, so uniqueness for~\eqref{eq:rivoal} proves the formulas.  This finite
calculation, including the two telescoping certificates and their boundary
terms, is checked symbolically and in Lean in
\texttt{RamanujanChallenge/Problem22.lean}.
\end{proof}

Let $\widetilde A_n=(n!)^2A_n$ and $\widetilde B_n=(n!)^2B_n$.  Define the
first-order operator
\begin{equation}\label{eq:ore}
 \mathcal O(T)_n=
 \frac{T_{n+1}+(n+1)(3n+4)T_n}{8n+11}.
\end{equation}

\begin{proposition}[exact Ore transform]\label{prop:ore}
The operator~\eqref{eq:ore} sends every factorial-scaled solution of
\eqref{eq:rivoal} to a solution of~\eqref{eq:challenge}, and
\[
 P_n=\mathcal O(\widetilde A)_n,qquad
 Q_n=\mathcal O(\widetilde B)_n.
\]
Moreover $Q_n>0$ for all $n$.
\end{proposition}

\begin{proof}
Insert~\eqref{eq:ore} into the left side of~\eqref{eq:challenge}.  After
clearing denominators, it is a polynomial linear combination of
\eqref{eq:rivoal} at $n$ and at $n+1$.  The transformed initial triples are
respectively $(0,7,179)$ and $(1,12,306)$, proving the identities by
uniqueness.  Positivity follows from~\eqref{eq:Bsum}: both terms in the
numerator of~\eqref{eq:ore} are positive.  The polynomial identity is also
machine-checked in the Lean source cited above.
\end{proof}

It remains to prove $A_n/B_n\to\gamma$.  Indeed, if $a_n=(n+1)(3n+4)$, then
\[
 \frac{\mathcal O(\widetilde A)_n}{\mathcal O(\widetilde B)_n}
 =\frac{\widetilde A_{n+1}+a_n\widetilde A_n}
        {\widetilde B_{n+1}+a_n\widetilde B_n}.
\]
This is a positive weighted average of
$\widetilde A_{n+1}/\widetilde B_{n+1}$ and
$\widetilde A_n/\widetilde B_n$, so it has the same limit.

\section{A finite Stein identity}

Use the positive weights
\[
 W_{n,k}=(2n+k+1)w_{n,k},\qquad
 Z_n=\sum_{k=0}^nW_{n,k}=B_n,
\]
and let $K$ be the random variable on $\{0,\ldots,n\}$ with
$\Pr(K=k)=W_{n,k}/Z_n$.  Define
\begin{align*}
 b_n(k)&=(2n+k+2)(n-k)^2,\\
 d_n(k)&=(2n+k)k^3,\\
 G_n(k)&=k^3-(n-k)^2.
\end{align*}
The adjacent-weight identity is exactly
\begin{equation}\label{eq:detailedbalance}
 W_{n,k}b_n(k)=W_{n,k+1}d_n(k+1).
\end{equation}
Since $b_n(n)=d_n(0)=0$, reindexing the finite sum gives, for every function
$f$,
\begin{equation}\label{eq:stein}
 \E\!\left[b_n(K)(f(K+1)-f(K))
 +(b_n(K)-d_n(K))f(K)\right]=0.
\end{equation}
Two elementary polynomial identities drive the estimate:
\begin{align}
G_n(k+1)-G_n(k)&=3k^2+k+2n,                                      \label{eq:deltaG}\\
b_n(k)-d_n(k)&=2(n-k)^2-(2n+k)G_n(k).                            \label{eq:bminusd}
\end{align}

\begin{lemma}[moment bounds]\label{lem:moments}
For $n\ge1$,
\begin{align}
 \E[K^3]&\le \frac52n^2,                                        \label{eq:k3}\\
 \E[K^2]&\le \frac72n\sqrt n,                                  \label{eq:k2}\\
 \E[G_n(K)^2]&\le81n^3\sqrt n.                                 \label{eq:G2}
\end{align}
\end{lemma}

\begin{proof}
Taking $f=1$ in~\eqref{eq:stein} gives $\E[b_n(K)]=\E[d_n(K)]$.
On the support,
\[
 d_n(k)\ge2nk^3,qquad b_n(k)\le5n^3,
\]
which proves~\eqref{eq:k3}.  The interpolation inequality
\[
 x^2\le n+\frac{x^3}{\sqrt n}\qquad(x\ge0)
\]
then gives~\eqref{eq:k2}.

Next take $f=G_n$ in~\eqref{eq:stein} and use
\eqref{eq:deltaG}--\eqref{eq:bminusd}.  We obtain the exact identity
\begin{align*}
 \E[(2n+K)G_n(K)^2]
  =\E\!\left[2(n-K)^2G_n(K)
   +b_n(K)(3K^2+K+2n)\right].
\end{align*}
Young's inequality and the support bounds give
\begin{align*}
2(n-k)^2G_n(k)&\le nG_n(k)^2+n^3,\\
b_n(k)(3k^2+k+2n)&\le20n^3k^2+10n^4.
\end{align*}
The left side is at least $2n\E[G_n(K)^2]$.  Hence
\[
 n\E[G_n(K)^2]
 \le20n^3\E[K^2]+11n^4
 \le81n^4\sqrt n,
\]
which proves~\eqref{eq:G2}.
\end{proof}

\section{From the cubic saddle to harmonic numbers}

The quotient in Proposition~\ref{prop:sums} is
\begin{equation}\label{eq:weightedaverage}
 \frac{A_n}{B_n}
 =\E\!\left[3H_K-2H_{n-K}-\frac1{2n+K+1}\right].
\end{equation}
We first omit the last, uniformly negligible, term.

For $m\ge1$, the standard elementary bounds for harmonic numbers imply
\begin{equation}\label{eq:harmonicremainder}
 0\le H_m-\log m-\gamma\le\frac1m.
\end{equation}
Fix $0<\delta\le1/16$ and call $k$ good when
\[
 |G_n(k)|\le\delta n^2.
\]
If $n\ge16$ and $k$ is good, direct comparison in
$k^3=(n-k)^2+G_n(k)$ gives
\begin{equation}\label{eq:geometry}
 1\le k<n,qquad k^3\ge\frac3{16}n^2,qquad n-k\ge\frac n2.
\end{equation}
Suppose also that $M\ge1$ and $16M^3\le3n^2$.  Then~\eqref{eq:geometry}
implies $k\ge M$ and $n-k\ge M$.  Put
\[
 x=\frac{k^3}{(n-k)^2}.
\]
The last inequality in~\eqref{eq:geometry} gives
\[
 |x-1|=\frac{|G_n(k)|}{(n-k)^2}\le4\delta\le\frac14.
\]
Since $|\log x|\le2|x-1|$ in this range,~\eqref{eq:harmonicremainder}
shows that every good $k$ satisfies
\begin{equation}\label{eq:goodbound}
 |3H_k-2H_{n-k}-\gamma|\le\frac5M+8\delta.
\end{equation}

On the bad set we use only the coarse estimate
\begin{equation}\label{eq:envelope}
 |3H_k-2H_{n-k}-\gamma|\le6+5\log n
 \qquad(0\le k\le n),
\end{equation}
which follows from $0\le H_m\le1+\log n$ and $0<\gamma<1$.
Badness also implies
\[
 1\le\frac{G_n(k)^2}{\delta^2n^4}.
\]
Combining~\eqref{eq:G2}, \eqref{eq:goodbound}, and
\eqref{eq:envelope} yields the quantitative estimate
\begin{equation}\label{eq:mainbound}
 \E\bigl[|3H_K-2H_{n-K}-\gamma|\bigr]
 \le \frac5M+8\delta+
 \frac{81(6+5\log n)}{\delta^2\sqrt n}.
\end{equation}

Given $\varepsilon>0$, first choose
$0<\delta\le\min(1/16,\varepsilon/32)$, then choose
$M>20/\varepsilon$.  For all sufficiently large $n$ the condition
$16M^3\le3n^2$ holds and the last term of~\eqref{eq:mainbound} is less
than $\varepsilon/2$.  Therefore the left side tends to zero.  Finally,
\[
 0<\E\!\left[\frac1{2n+K+1}\right]\le\frac1{2n+1}\longrightarrow0.
\]
Equation~\eqref{eq:weightedaverage} now proves
\begin{equation}\label{eq:rivoallimit}
 \frac{A_n}{B_n}\longrightarrow\gamma.
\end{equation}

\section{Conclusion}

By Proposition~\ref{prop:ore}, the challenge quotient is a positive weighted
average of two adjacent quotients in~\eqref{eq:rivoallimit}.  Consequently
\[
 \boxed{\displaystyle\lim_{n\to\infty}\frac{p_n}{q_n}=\gamma}.
\]

The complete argument has been formalized in Lean~4, in namespace
\texttt{RamanujanChallenge.P22}.  The public theorem is
\texttt{problem22\_solved}.  Its axiom audit contains only the three standard
Lean foundations \texttt{propext}, \texttt{Classical.choice}, and
\texttt{Quot.sound}.  The proof contains no \texttt{sorry} declarations or
problem-specific axioms.

\end{document}
